\begin{figure}[t]
\centering
\begin{tikzpicture}[
  font=\footnotesize,
  box/.style={draw, rounded corners=2pt, align=center,
              minimum height=8mm, inner sep=2.5pt},
  arrow/.style={-{Latex[length=1.8mm]}, thick},
  note/.style={font=\scriptsize, align=center}
]
\node[box, fill=stagegray, text width=.78\columnwidth] (screen)
  {\textbf{동일한 웹 화면과 Stage 1 evidence}\\
   content · state · relation · affordance · coordinate};
\node[box, fill=shortblue, below=7mm of screen,
      xshift=-.27\columnwidth, text width=.25\columnwidth] (concise)
  {\textbf{Concise}\\salient evidence};
\node[box, fill=stagegreen, below=7mm of screen,
      text width=.25\columnwidth] (balanced)
  {\textbf{Balanced}\\section-level\\evidence};
\node[box, fill=longorange, below=7mm of screen,
      xshift=.27\columnwidth, text width=.25\columnwidth] (comprehensive)
  {\textbf{Compre-}\\\textbf{hensive}\\near-complete\\evidence};
\node[box, fill=stagegray, below=10mm of balanced,
      text width=.78\columnwidth] (eval)
  {\textbf{공통 grounding 및 후속 학습}\\
   coverage 분석 · targeted diagnostic\\공개 benchmark · agent transfer};
\draw[arrow] (screen.south) -- (concise.north);
\draw[arrow] (screen.south) -- (balanced.north);
\draw[arrow] (screen.south) -- (comprehensive.north);
\draw[arrow] (concise.south) -- (eval.north west);
\draw[arrow] (balanced.south) -- (eval.north);
\draw[arrow] (comprehensive.south) -- (eval.north east);
\end{tikzpicture}
\caption{\textbf{웹페이지 verbalization의 정보 포괄성 비교.}
동일한 화면과 Stage 1 evidence를 서로 다른 information budget으로
언어화한다. 세 recipe는 평균 coverage가 증가하지만 sample별
evidence 집합의 엄격한 포함관계를 가정하지 않는다. Grounding
데이터와 후속 학습은 동일하게 유지한다.}
\label{fig:contrast}
\end{figure}
